\chapter{Project Overiew}

This project aims to address a critical gap in multi-view Scene Change Detection (SCD) for underwater environments by investigating and developing primitive-space change detection methods within underwater 3D Gaussian Splatting (3DGS) models. 

The primary objective is to move beyond traditional render-then-compare paradigms (which rely on pixel-space comparison) by performing change detection directly on the 3DGS primitives. By analyzing native primitive attributes—position, anisotropic covariance, and color—this project will evaluate whether these attributes carry sufficient signal for reliable scene change detection. 

The research will investigate algorithms to handle the inherent challenges of independent 3DGS reconstructions, specifically geometric drift (representation ambiguity and observation uncertainty) and photometric drift. Ultimately, the project will apply this methodology to the EcoRRAP dataset provided by the Australian Institute of Marine Science (AIMS) and the QUT Centre for Robotics, testing its ability to separate structural changes (e.g., coral growth or breakage) from surface-level appearance changes (e.g., bleaching or lighting shifts) in a complex, real-world underwater environment.
